\documentclass[10pt,a4paper]{amsart}
\usepackage{amsmath,amssymb,amsthm,graphicx}
% Retained preamble for the Stacks Project chapter snapshot in this directory.
% The chapter's own preamble.tex is a copy of the shared Stacks Project preamble; the block
% below is sliced MECHANICALLY, line for line, from it. Nothing was added, removed or
% redefined; the wrapper class, the xr-hyper/hyperref wiring and the externaldocument list
% are omitted because the assembler supplies the document class.
%
% lines taken from preamble.tex: 3, 5, 9, 11, 12, 156, 157, 158, ... (51 lines total)
%
% package coverage restored from the chapter's own preamble (2026-10-13): xy, xy-2cell

\usepackage[all]{xy}
\xyoption{2cell}
\UseAllTwocells

\newenvironment{reference}{\comment}{\endcomment}
\newenvironment{slogan}{\comment}{\endcomment}
\newenvironment{history}{\comment}{\endcomment}
\theoremstyle{plain}
\newtheorem{theorem}[subsection]{Theorem}
\newtheorem{proposition}[subsection]{Proposition}
\newtheorem{lemma}[subsection]{Lemma}
\theoremstyle{definition}
\newtheorem{definition}[subsection]{Definition}
\newtheorem{example}[subsection]{Example}
\newtheorem{exercise}[subsection]{Exercise}
\newtheorem{situation}[subsection]{Situation}
\theoremstyle{remark}
\newtheorem{remark}[subsection]{Remark}
\newtheorem{remarks}[subsection]{Remarks}
\numberwithin{equation}{subsection}
\def\lim{\mathop{\mathrm{lim}}\nolimits}
\def\colim{\mathop{\mathrm{colim}}\nolimits}
\def\Spec{\mathop{\mathrm{Spec}}}
\def\Hom{\mathop{\mathrm{Hom}}\nolimits}
\def\Ext{\mathop{\mathrm{Ext}}\nolimits}
\def\SheafHom{\mathop{\mathcal{H}\!\mathit{om}}\nolimits}
\def\SheafExt{\mathop{\mathcal{E}\!\mathit{xt}}\nolimits}
\def\Sch{\mathit{Sch}}
\def\Mor{\mathop{\mathrm{Mor}}\nolimits}
\def\Ob{\mathop{\mathrm{Ob}}\nolimits}
\def\Sh{\mathop{\mathit{Sh}}\nolimits}
\def\NL{\mathop{N\!L}\nolimits}
\def\CH{\mathop{\mathrm{CH}}\nolimits}
\def\proetale{{pro\text{-}\acute{e}tale}}
\def\etale{{\acute{e}tale}}
\def\QCoh{\mathit{QCoh}}
\def\Ker{\mathop{\mathrm{Ker}}}
\def\Im{\mathop{\mathrm{Im}}}
\def\Coker{\mathop{\mathrm{Coker}}}
\def\Coim{\mathop{\mathrm{Coim}}}
\DeclareMathSymbol{\boxtimes}{\mathbin}{AMSa}{"02}
\def\QCohstack{\mathcal{QC}\!\mathit{oh}}
\def\Cohstack{\mathcal{C}\!\mathit{oh}}
\def\Spacesstack{\mathcal{S}\!\mathit{paces}}
\def\Quotfunctor{\mathrm{Quot}}
\def\Hilbfunctor{\mathrm{Hilb}}
\def\Curvesstack{\mathcal{C}\!\mathit{urves}}
\def\Polarizedstack{\mathcal{P}\!\mathit{olarized}}
\def\Complexesstack{\mathcal{C}\!\mathit{omplexes}}
\def\Pic{\mathop{\mathrm{Pic}}\nolimits}
\def\Picardstack{\mathcal{P}\!\mathit{ic}}
\def\Picardfunctor{\mathrm{Pic}}
\def\Deformationcategory{\mathcal{D}\!\mathit{ef}}
\usepackage{cleveref}
\crefname{originallemma}{original lemma}{original lemmas}
\Crefname{originallemma}{Original lemma}{Original lemmas}
\makeatletter
\newlabel{lemma-space-presentation}{{0262}{source}{}{sourcepointer-1}{}}
\newlabel{lemma-space-presentation@cref}{{[originallemma][0][]0262}{[1][0][]source}{}{}{}}
\makeatother
\crefname{originallemma}{original lemma}{original lemmas}
\Crefname{originallemma}{Original lemma}{Original lemmas}
\makeatletter
\newlabel{sites-lemma-describe-j-shriek-good-site}{{03HT}{source}{}{sourcepointer-2}{}}
\newlabel{sites-lemma-describe-j-shriek-good-site@cref}{{[originallemma][0][]03HT}{[1][0][]source}{}{}{}}
\makeatother
\crefname{originallemma}{original lemma}{original lemmas}
\Crefname{originallemma}{Original lemma}{Original lemmas}
\makeatletter
\newlabel{sites-lemma-describe-j-shriek-representable}{{03HU}{source}{}{sourcepointer-3}{}}
\newlabel{sites-lemma-describe-j-shriek-representable@cref}{{[originallemma][0][]03HU}{[1][0][]source}{}{}{}}
\makeatother
\crefname{originallemma}{original lemma}{original lemmas}
\Crefname{originallemma}{Original lemma}{Original lemmas}
\makeatletter
\newlabel{sites-lemma-j-shriek-commutes-equalizers-fibre-products}{{04BB}{source}{}{sourcepointer-4}{}}
\newlabel{sites-lemma-j-shriek-commutes-equalizers-fibre-products@cref}{{[originallemma][0][]04BB}{[1][0][]source}{}{}{}}
\makeatother
\crefname{originallemma}{original lemma}{original lemmas}
\Crefname{originallemma}{Original lemma}{Original lemmas}
\makeatletter
\newlabel{sites-lemma-localize-given-products}{{03CE}{source}{}{sourcepointer-5}{}}
\newlabel{sites-lemma-localize-given-products@cref}{{[originallemma][0][]03CE}{[1][0][]source}{}{}{}}
\makeatother
\crefname{originallemma}{original lemma}{original lemmas}
\Crefname{originallemma}{Original lemma}{Original lemmas}
\makeatletter
\newlabel{sites-lemma-relocalize}{{03EH}{source}{}{sourcepointer-6}{}}
\newlabel{sites-lemma-relocalize@cref}{{[originallemma][0][]03EH}{[1][0][]source}{}{}{}}
\makeatother
\crefname{originalsection}{original section}{original sections}
\Crefname{originalsection}{Original section}{Original sections}
\makeatletter
\newlabel{sites-section-localize}{{00XZ}{source}{}{sourcepointer-7}{}}
\newlabel{sites-section-localize@cref}{{[originalsection][0][]00XZ}{[1][0][]source}{}{}{}}
\makeatother
\crefname{originaltheorem}{original theorem}{original theorems}
\Crefname{originaltheorem}{Original theorem}{Original theorems}
\makeatletter
\newlabel{theorem-presentation}{{02WW}{source}{}{sourcepointer-8}{}}
\newlabel{theorem-presentation@cref}{{[originaltheorem][0][]02WW}{[1][0][]source}{}{}{}}
\makeatother
\title{Localization functors preserve schemes and algebraic spaces under a change of base}
\author{The Stacks Project Authors}
\date{stacks-project:spaces.tex:local-snapshot-20261012; GNU FDL 1.2 or later}
\begin{document}
\maketitle
\setcounter{subsection}{1}
\noindent\textit{Editorial note.} All delivered mathematical text is the author's own, copied from the pinned Stacks Project chapter file spaces.tex; no statement or proof step was rewritten, reordered or removed.\par\medskip
\noindent\textit{Editorial note.} The statement is a six-part assertion; the two enumerate blocks are part of the counted statement span and are preserved verbatim.\par\medskip
\noindent\textit{Editorial note.} The eight references of the counted spans are recorded as non-load-bearing navigation pointers; the argument that consumes them stays inside the counted proof span.\par\medskip
\noindent\textit{Editorial note.} Preamble: sources/stacks-spaces/retained-preamble.tex, sliced mechanically from the chapter's own preamble.tex.\par\medskip
\noindent\textit{Allowed standard prerequisites (editorial).}\begin{itemize}
\item Definition of the localization functors j\_!, j\_*, j\textasciicircum{}\{-1\}. Sites, Lemma sites-lemma-relocalize, used at proof lines 2982-2986. Makes the localization material available for j.
\item The coproduct formula for j\_!. Sites, Lemma sites-lemma-describe-j-shriek-good-site, used at proof lines 2987-2988. Gives assertion (1).
\item j\_! commutes with fibre products and equalizers. Sites, Lemma sites-lemma-j-shriek-commutes-equalizers-fibre-products, used at proof lines 2990-2992. Transports the equivalence relation.
\item j\_! of a representable sheaf is representable. Sites, Lemma sites-lemma-describe-j-shriek-representable, used at proof lines 2993-2995. Gives assertion (2).
\item Presentation of an algebraic space and the presentation theorem. Lemma lemma-space-presentation (1387-1407) and Theorem theorem-presentation, used at proof lines 2996-3009 and 3018-3024. Recognizes algebraic spaces from etale equivalence relations.
\end{itemize}\medskip
\ifdefined\hypertarget\hypertarget{sourcepointer-1}{}\fi
\noindent\textit{Source reference.} Original lemma 0262: focused-seeds/AGP/sources/stacks-snapshot-20261011/tags.txt, source lines 2835--2835. Navigation pointer into the Stacks Project; the standard fact it supplies is applied inside the counted proof span and the argument resting on it stays there.\par\medskip
\ifdefined\hypertarget\hypertarget{sourcepointer-2}{}\fi
\noindent\textit{Source reference.} Original lemma 03HT: focused-seeds/AGP/sources/stacks-snapshot-20261011/tags.txt, source lines 4579--4579. Navigation pointer into the Stacks Project; the standard fact it supplies is applied inside the counted proof span and the argument resting on it stays there.\par\medskip
\ifdefined\hypertarget\hypertarget{sourcepointer-3}{}\fi
\noindent\textit{Source reference.} Original lemma 03HU: focused-seeds/AGP/sources/stacks-snapshot-20261011/tags.txt, source lines 4580--4580. Navigation pointer into the Stacks Project; the standard fact it supplies is applied inside the counted proof span and the argument resting on it stays there.\par\medskip
\ifdefined\hypertarget\hypertarget{sourcepointer-4}{}\fi
\noindent\textit{Source reference.} Original lemma 04BB: focused-seeds/AGP/sources/stacks-snapshot-20261011/tags.txt, source lines 5646--5646. Navigation pointer into the Stacks Project; the standard fact it supplies is applied inside the counted proof span and the argument resting on it stays there.\par\medskip
\ifdefined\hypertarget\hypertarget{sourcepointer-5}{}\fi
\noindent\textit{Source reference.} Original lemma 03CE: focused-seeds/AGP/sources/stacks-snapshot-20261011/tags.txt, source lines 4381--4381. Navigation pointer into the Stacks Project; the standard fact it supplies is applied inside the counted proof span and the argument resting on it stays there.\par\medskip
\ifdefined\hypertarget\hypertarget{sourcepointer-6}{}\fi
\noindent\textit{Source reference.} Original lemma 03EH: focused-seeds/AGP/sources/stacks-snapshot-20261011/tags.txt, source lines 4456--4456. Navigation pointer into the Stacks Project; the standard fact it supplies is applied inside the counted proof span and the argument resting on it stays there.\par\medskip
\ifdefined\hypertarget\hypertarget{sourcepointer-7}{}\fi
\noindent\textit{Source reference.} Original section 00XZ: focused-seeds/AGP/sources/stacks-snapshot-20261011/tags.txt, source lines 1237--1237. Navigation pointer into the Stacks Project; the standard fact it supplies is applied inside the counted proof span and the argument resting on it stays there.\par\medskip
\ifdefined\hypertarget\hypertarget{sourcepointer-8}{}\fi
\noindent\textit{Source reference.} Original theorem 02WW: focused-seeds/AGP/sources/stacks-snapshot-20261011/tags.txt, source lines 3819--3819. Navigation pointer into the Stacks Project; the standard fact it supplies is applied inside the counted proof span and the argument resting on it stays there.\par\medskip
\begin{lemma}
\label{lemma-change-base-scheme}
Suppose given a big site $\Sch_{fppf}$.
Let $g : S \to S'$ be morphism of $\Sch_{fppf}$.
Let $j : (\Sch/S)_{fppf} \to (\Sch/S')_{fppf}$ be
the corresponding localization functor.
Let $F$ be a sheaf of sets on $(\Sch/S)_{fppf}$.
Then
\begin{enumerate}
\item for a scheme $T'$ over $S'$ we have
$j_!F(T'/S') =
\coprod\nolimits_{\varphi : T' \to S} F(T' \xrightarrow{\varphi} S),$
\item if $F$ is representable by a scheme
$X \in \Ob((\Sch/S)_{fppf})$,
then $j_!F$ is representable by $j(X)$ which is
$X$ viewed as a scheme over $S'$, and
\item if $F$ is an algebraic space over $S$, then $j_!F$ is an algebraic
space over $S'$, and if $F = U/R$ is a presentation, then
$j_!F = j(U)/j(R)$ is a presentation.
\end{enumerate}
Let $F'$ be a sheaf of sets on $(\Sch/S')_{fppf}$. Then
\begin{enumerate}
\item[(4)] for a scheme $T$ over $S$ we have $j^{-1}F'(T/S) = F'(T/S')$,
\item[(5)] if $F'$ is representable by a scheme
$X' \in \Ob((\Sch/S')_{fppf})$, then
$j^{-1}F'$ is representable, namely by $X'_S = S \times_{S'} X'$, and
\item[(6)] if $F'$ is an algebraic space, then
$j^{-1}F'$ is an algebraic space, and if $F' = U'/R'$ is a presentation,
then $j^{-1}F' = U'_S/R'_S$ is a presentation.
\end{enumerate}
\end{lemma}

\begin{proof}
The functors $j_!$, $j_*$ and $j^{-1}$ are defined in
Sites, Lemma \ref{sites-lemma-relocalize}
where it is also shown that $j = j_{S/S'}$ is the localization
of $(\Sch/S')_{fppf}$ at the object $S/S'$. Hence
all of the material on localization functors is available for $j$.
The formula in (1) is
Sites, Lemma \ref{sites-lemma-describe-j-shriek-good-site}.
By definition $j_!$ is the left adjoint to restriction $j^{-1}$,
hence $j_!$ is right exact. By
Sites, Lemma \ref{sites-lemma-j-shriek-commutes-equalizers-fibre-products}
it also commutes with fibre products and equalizers.
By
Sites, Lemma \ref{sites-lemma-describe-j-shriek-representable}
we see that $j_!h_X = h_{j(X)}$ hence (2) holds.
If $F$ is an algebraic space over $S$, then we can write $F = U/R$
(Lemma \ref{lemma-space-presentation})
and we get
$$
j_!F = j(U)/j(R)
$$
because $j_!$ being right exact commutes with coequalizers, and moreover
$j(R) = j(U) \times_{j_!F} j(U)$ as $j_!$ commutes with fibre products.
Since the morphisms $j(s), j(t) : j(R) \to j(U)$ are simply the morphisms
$s, t : R \to U$ (but viewed as morphisms of schemes over $S'$), they
are still \'etale. Thus $(j(U), j(R), s, t)$ is an \'etale equivalence relation.
Hence by
Theorem \ref{theorem-presentation}
we conclude that $j_!F$ is an algebraic space.

\medskip\noindent
Proof of (4), (5), and (6). The description of $j^{-1}$ is in
Sites, Section \ref{sites-section-localize}.
The restriction of the representable sheaf associated to $X'/S'$
is the representable sheaf associated to
$X'_S = S \times_{S'} Y'$ by
Sites, Lemma \ref{sites-lemma-localize-given-products}.
The restriction functor $j^{-1}$ is exact, hence $j^{-1}F' = U'_S/R'_S$.
Again by exactness the sheaf $R'_S$ is still an equivalence relation on
$U'_S$. Finally the two maps $R'_S \to U'_S$ are \'etale as base changes
of the \'etale morphisms $R' \to U'$. Hence $j^{-1}F' = U'_S/R'_S$ is
an algebraic space by
Theorem \ref{theorem-presentation}
and we win.
\end{proof}

\end{document}
